% !PRINTED 639
\weberpage{639}
%<<CONT
sumptions, from which we shall afterwards free the result again.
These assumptions are:

\gesperrt{1. that by this substitution of the $\xi$ for the $x$ none of
the resolvents $(\varepsilon, \xi)$, in which $\varepsilon$ is an
imaginary $n$th root of unity, vanishes.}

By (13), (16), (17), §.184 the functions $f_\nu, F_\nu,
\Phi_\nu$ then receive definite values different from zero, and we make
the second assumption,

% [sic] Print runs the list through $f_{n-1}$, although the system defined in
% §.184 (5) and used immediately below runs only from $f_0$ through $f_{n-2}$.
\gesperrt{2. that by the same substitution no two of the functions
$f_0, f_1 \ldots f_{n-1}$ become equal to one another.}

We assume that by the substitution of the $\xi$ for the $x$

\[
f_0, f_1 \ldots f_{n-2}
\]

pass into

\[
k_0, k_1 \ldots k_{n-2},
\]

so that the $k_0, k_1 \ldots k_{n-2}$ are therefore different from one
another, and

\[
\Phi_0, \Phi_1 \ldots \Phi_{n-2}
\]

into

\[
K_0, K_1 \ldots K_{n-2}
\]

where, however, among the $K$ equal ones may also occur.

On account of the property $\delta)$ of the functions $f_\nu$, the
quantities $k_0, k_1 \ldots k_{n-2}$ are the roots of a cyclic equation
of degree $(n-1)$ (hence even), $\psi(u) = 0$, so that, if $u$ denotes a
variable,

\weq{1}{\psi(u) = (u - k_0)(u - k_1) \cdots (u - k_{n-2})}

is a function of degree $(n-1)$ in $u$ with rational coefficients. The
quantities $k_0, k_1 \ldots k_{n-2}$ can be expressed rationally in
terms of one another in the form

\weq{2}{k_1 = \Theta(k_0), \quad k_2 = \Theta(k_1), \ldots k_{n-1} = \Theta(k_{n-2}), \quad k_0 = \Theta(k_{n-1})} % [sic] Print refers to $k_{n-1}$ although the sequence is introduced only through $k_{n-2}$; preserved.
(§.\,163).

Likewise the quantities $K_\nu$ can, by the concluding theorem of
§.184, be expressed in terms of the $k_\nu$, and indeed in the form

\weq{3}{K_0 = \Phi(k_0), \quad K_1 = \Phi(k_1), \ldots K_{n-2} = \Phi(k_{n-2}),}

in which $\Phi$ denotes a rational function (in $\mathfrak{K}$).

Accordingly the formula (18), §.184 yields us the following result: We
put, as an abbreviation,

\weq{4}{\tau_\nu = \sqrt[n]{k_\nu},}
%CONT>>
